\section{Evaluation Methodology}
\label{sec:evaluation}

\subsection{Metrics}

Model performance is evaluated using four quantitative metrics, computed on images normalized to $[0, 1]$.

\subsubsection{PSNR (Peak Signal-to-Noise Ratio)}
\begin{equation}
\text{PSNR} = 20 \cdot \log_{10}\!\left(\frac{\text{MAX}}{\sqrt{\text{MSE}}}\right) \quad [\text{dB}]
\end{equation}
where $\text{MAX} = 1.0$. Higher values indicate greater fidelity to ground truth.

\subsubsection{SSIM (Structural Similarity Index)}
SSIM~\cite{ssim} is computed with an $11\!\times\!11$ window via \texttt{F.avg\_pool2d}:
\begin{equation}
\text{SSIM}(x, y) = \frac{(2\mu_x\mu_y + C_1)(2\sigma_{xy} + C_2)}{(\mu_x^2 + \mu_y^2 + C_1)(\sigma_x^2 + \sigma_y^2 + C_2)}
\end{equation}
with $C_1 = (0.01)^2$ and $C_2 = (0.03)^2$. Values closer to 1 indicate higher structural similarity.

\subsubsection{MAE and MSE}
Mean Absolute Error and Mean Squared Error, computed pixel-wise via \texttt{F.l1\_loss} and \texttt{F.mse\_loss} respectively.

\subsection{Full-Resolution Inference}

DIV2K images have resolution ${\sim}2000\!\times\!2000$ pixels, too large for a single forward pass. Inference is performed via \textbf{sliding window} with weighted blending.

\subsubsection{Sliding Window}
The image is partitioned into overlapping patches of size \texttt{patch\_size} (typically 256) with \texttt{overlap} pixels (typically 32--64) of overlap:
\begin{equation}
\text{stride} = \text{patch\_size} - \text{overlap}
\end{equation}

\subsubsection{Hann Window Blending}
To avoid boundary artifacts between patches, each output is weighted with a 2D Hann window, obtained as the outer product of two 1D windows:
\begin{equation}
w(i, j) = w_H(i) \cdot w_H(j), \quad w_H(k) = 0.5 - 0.5 \cos\!\left(\frac{2\pi k}{N - 1}\right)
\end{equation}

The final output for each pixel is the weighted average of contributions from all overlapping patches:
\begin{equation}
I_{\text{out}}(i, j) = \frac{\sum_p w_p \cdot O_p(i, j)}{\sum_p w_p}
\end{equation}
This produces smooth, natural transitions between adjacent patches.

\subsection{Evaluation Protocol}

The \texttt{ImageRestorationEvaluator} class orchestrates evaluation on the full dataset:
\begin{enumerate}
    \item For each (degraded, clean) pair: sliding-window inference followed by per-image metric computation.
    \item Aggregation: mean and standard deviation across all images.
    \item Reporting: JSON export with per-image and aggregate metrics, identification of best and worst cases.
\end{enumerate}

All models are evaluated on the same 10 validation images (indices 0801--0810 from DIV2K), enabling direct per-image comparisons across architectures.
